\documentclass[12pt,a4paper]{article}
\usepackage[utf8]{inputenc}
\usepackage[T1]{fontenc}
\usepackage{times}
\usepackage[a4paper,top=2.5cm,bottom=2.5cm,left=3cm,right=2.5cm]{geometry}
\usepackage{fancyvrb}
\usepackage{enumitem}
\setlength{\parindent}{0pt}
\setlength{\parskip}{6pt}
\setlist{itemsep=1pt,topsep=3pt}
\renewcommand{\refname}{Referências}
\begin{document}

\begin{center}
{\bfseries Pontifícia Universidade Católica de Minas Gerais}\\
Instituto de Ciências Exatas e Informática (ICEI)\\
Engenharia de Software\\
Disciplina: Sistemas Operacionais\\
Professor: Diego Silva Caldeira Rocha\\[12pt]
{\Large\bfseries Trabalho Prático 1}\\
{\large Comunicação e Sincronização de Processos e Threads}\\[12pt]
\end{center}

\textbf{Integrantes:}\\
Lucas Erick Andrade da Silva\\
Julia Carvalho\\
João Paulo de Freitas Fadul\\
Victhor Guilherme Santos Loiola\\
Bernardo Demaria Santos

\section*{1. Introdução}
Neste trabalho nós estudamos como várias threads podem usar o mesmo recurso sem causar erro. Na Parte I, rodamos o
exemplo do produtor e consumidor que o professor passou. Na Parte II, fizemos em Java os sete problemas do
enunciado, usando threads e semáforos. Cada programa tem um teste que confere se tudo funcionou direito.

Algumas palavras que aparecem no texto:
\begin{itemize}
  \item \textbf{Thread:} uma tarefa do programa que roda ao mesmo tempo que outras.
  \item \textbf{Semáforo:} um contador que controla quem pode continuar. Se o valor é zero, a thread espera.
  \item \textbf{Região crítica:} o pedaço do código que só uma thread pode executar por vez.
  \item \textbf{Deadlock:} quando as threads ficam esperando umas pelas outras e o programa trava.
\end{itemize}

\section*{2. Parte I: Thread em Java}
A classe \texttt{Buffer} guarda um valor só e uma variável \texttt{cheio}. O método \texttt{put} só grava quando
o buffer está vazio, e o método \texttt{get} só lê quando ele está cheio. Os dois são \texttt{synchronized}, ou
seja, só uma thread entra por vez. Quando não pode continuar, a thread chama \texttt{wait()} e fica esperando.
Quem termina chama \texttt{notifyAll()} para acordar a outra.

O \texttt{Produtor} coloca os valores de 0 a 4 e o \texttt{Consumidor} retira 5 valores. Como o buffer só guarda
um valor, um novo \texttt{put} só acontece depois do \texttt{get} do valor anterior.

Esta é a saída de uma execução. As linhas \texttt{put} e \texttt{get} são impressas dentro do buffer, e as linhas
\texttt{Produtor} e \texttt{Consumidor} são impressas fora dele. Por isso a ordem entre elas pode mudar de uma
execução para outra, mas os valores sempre saem em ordem (0, 1, 2, 3, 4) e nenhum se perde.
\begin{small}
\begin{Verbatim}[frame=single]
put 0
get 0
Produtor #0 put: 0
Consumidor #0 get: 0
put 1
Produtor #0 put: 1
get 1
Consumidor #0 get: 1
put 2
get 2
Consumidor #0 get: 2
Produtor #0 put: 2
put 3
get 3
Consumidor #0 get: 3
Produtor #0 put: 3
put 4
get 4
Consumidor #0 get: 4
Produtor #0 put: 4
\end{Verbatim}
\end{small}

\section*{3. Parte II: Programas com semáforos}
Os programas estão na pasta \texttt{src/trabalho}. Em cada problema abaixo dizemos quais são as threads, o que elas
dividem, qual é a região crítica, quais semáforos usamos (com o valor inicial), como evitamos deadlock e como
testamos. Se um programa travar, o teste acusa erro depois de um tempo, em vez de ficar parado para sempre.

\subsection*{3.1 Produtor e consumidor}
\textbf{Threads:} 3 produtores e 3 consumidores.\\
\textbf{O que dividem:} um buffer de 5 posições.\\
\textbf{Região crítica:} mexer no buffer e nas posições de entrada e saída.\\
\textbf{Semáforos:} \texttt{vazios} (começa em 5, conta as posições livres), \texttt{cheios} (começa em 0, conta os
itens guardados) e \texttt{mutex} (começa em 1, deixa só uma thread mexer no buffer).\\
\textbf{Deadlock:} a thread sempre espera \texttt{vazios} ou \texttt{cheios} antes de pegar o \texttt{mutex}, e solta
o \texttt{mutex} antes de avisar a outra. Assim ninguém dorme segurando o \texttt{mutex}.\\
\textbf{Teste:} os 60 itens produzidos são consumidos uma vez só, o buffer nunca passa de 5 itens e termina vazio.

\subsection*{3.2 Barbeiro dorminhoco}
\textbf{Threads:} 1 barbeiro e 12 clientes. A barbearia tem 3 cadeiras de espera.\\
\textbf{O que dividem:} a variável \texttt{waiting}, que conta os clientes esperando.\\
\textbf{Região crítica:} ler e mudar o \texttt{waiting}.\\
\textbf{Semáforos:} \texttt{customers} (0, conta os clientes esperando, e o barbeiro dorme nele), \texttt{barbers}
(0, avisa que o barbeiro está livre), \texttt{mutex} (1, protege o \texttt{waiting}) e \texttt{done} (0, avisa
que o corte terminou).\\
\textbf{Deadlock:} o cliente solta o \texttt{mutex} antes de esperar o barbeiro, e o barbeiro solta o \texttt{mutex}
antes de cortar. Se a barbearia está cheia, o cliente vai embora sem esperar em semáforo nenhum.\\
\textbf{Teste:} os atendidos mais os que foram embora somam 12, o barbeiro nunca corta dois ao mesmo tempo e a fila
nunca passa de 3.

\subsection*{3.3 Leitores e escritores}
\textbf{Threads:} 5 leitores e 2 escritores.\\
\textbf{O que dividem:} o vetor \texttt{M}, com 4 números.\\
\textbf{Região crítica:} escrever no vetor. Só um escritor por vez, e sem leitor junto. Já os leitores podem ler
ao mesmo tempo, porque ler não muda nada.\\
\textbf{Semáforos:} \texttt{mutex} (1, protege o contador de leitores), \texttt{db} (1, dá acesso exclusivo ao
vetor) e \texttt{turno} (1, é justo e funciona como uma catraca, para os leitores não deixarem os escritores
esperando para sempre).\\
\textbf{Deadlock:} o primeiro leitor pega o \texttt{db} e o último solta. O \texttt{mutex} é solto logo depois de
mexer no contador, e o \texttt{turno} é solto antes de ler ou escrever. Nenhum semáforo fica preso esperando outro.\\
\textbf{Teste:} o escritor grava o mesmo número nas 4 posições, devagar de propósito. Se um leitor visse números
diferentes, ou se um escritor rodasse junto com outra thread, o teste falharia.

\subsection*{3.4 Fumantes}
\textbf{Threads:} 1 representante e 3 fumantes (um com tabaco, um com papel e um com fósforos).\\
\textbf{O que dividem:} a mesa.\\
\textbf{Região crítica:} colocar e tirar os ingredientes da mesa.\\
\textbf{Semáforos:} \texttt{agente} (1, só deixa o representante repor depois que o fumante terminou) e um
semáforo para cada fumante, \texttt{fumante[i]} (0, o fumante fica parado até ser chamado).\\
\textbf{Deadlock:} só um fumante é chamado por vez, e ele sempre devolve o \texttt{agente} no fim. Por isso uma thread
nunca fica esperando outra que está esperando ela.\\
\textbf{Teste:} em 15 rodadas, o fumante chamado sempre achou na mesa os dois ingredientes que ele não tem, e foram
15 cigarros no total.

\subsection*{3.5 Filósofos}
\textbf{Threads:} 5 filósofos.\\
\textbf{O que dividem:} 5 palitos.\\
\textbf{Região crítica:} comer, que usa o palito da direita ($p_i$) e o da esquerda ($p_{i+1}$).\\
\textbf{Semáforos:} \texttt{palito[i]} (1 para cada palito) e \texttt{sala} (4).\\
\textbf{Deadlock:} se os 5 pegassem o palito da direita ao mesmo tempo, todos ficariam esperando o da esquerda e
o programa travaria. O semáforo \texttt{sala} deixa só 4 filósofos tentarem pegar os palitos. Assim sempre sobra um
palito livre e pelo menos um filósofo consegue comer.\\
\textbf{Teste do deadlock:} tiramos a \texttt{sala} (valor 5) e fizemos o filósofo esperar um pouco depois do primeiro
palito. O programa travou em 2 de 3 execuções. Com a \texttt{sala} em 4, todas terminaram.\\
\textbf{Teste:} dois vizinhos nunca comem juntos e cada filósofo come 6 vezes.

\subsection*{3.6 Jantar dos selvagens}
\textbf{Threads:} 4 selvagens e 1 cozinheiro.\\
\textbf{O que dividem:} o caldeirão, que tem 5 porções.\\
\textbf{Região crítica:} se servir e encher o caldeirão.\\
\textbf{Semáforos:} \texttt{mutex} (1, só um selvagem se serve por vez), \texttt{vazio} (0, acorda o cozinheiro) e
\texttt{cheio} (0, o cozinheiro avisa que encheu).\\
\textbf{Deadlock:} o selvagem que acha o caldeirão vazio acorda o cozinheiro e espera segurando o \texttt{mutex},
de propósito, para ninguém se servir de um caldeirão vazio. Não trava porque o cozinheiro não precisa do
\texttt{mutex} para encher.\\
\textbf{Teste:} nunca sobra porção negativa, o cozinheiro só enche quando está vazio e, com 24 refeições e 5 porções,
acontecem exatamente 4 recargas.

\subsection*{3.7 Três robôs}
\textbf{Threads:} Bart, Lisa e Maggie.\\
\textbf{O que dividem:} a vez de se mexer (\texttt{move()}).\\
\textbf{Semáforos:} \texttt{bart} (1, ele começa), \texttt{lisa} (0) e \texttt{maggie} (0). Só existe uma vez
circulando, então só um robô anda por vez e ninguém fica esperando à toa.\\
\textbf{Ordem:} Bart, Lisa, Maggie, Lisa, Bart, e assim por diante. A Lisa anda duas vezes por ciclo: no primeiro
movimento ela libera a Maggie e no segundo ela libera o Bart.\\
\textbf{Deadlock:} depois de se mexer, cada robô passa a vez para o próximo, então ninguém espera em círculo.\\
\textbf{Teste:} em 5 ciclos, os 20 movimentos seguem exatamente a ordem Bart, Lisa, Maggie, Lisa.

\section*{4. Testes}
Rodamos cada programa 20 vezes seguidas, porque erros de sincronização às vezes só aparecem em algumas execuções.
Todos passaram nas 20.
\begin{small}
\begin{Verbatim}[frame=single]
Produtor/Consumidor          20/20 execucoes OK
Barbeiro Dorminhoco          20/20 execucoes OK
Leitores e Escritores        20/20 execucoes OK
Fumantes                     20/20 execucoes OK
Filosofos                    20/20 execucoes OK
Selvagens                    20/20 execucoes OK
Tres Robos                   20/20 execucoes OK
TODOS OS TESTES PASSARAM
\end{Verbatim}
\end{small}
A saída completa de cada programa está na pasta \texttt{saidas}.

\section*{5. Como compilar e executar}
\begin{Verbatim}[frame=single]
javac -d out $(find src -name "*.java")
java -cp out ProdutorConsumidor.ProdCons     # Parte I
java -cp out trabalho.Main 1                 # problemas de 1 a 7
java -cp out trabalho.Main todos             # resumo de todos
java -cp out trabalho.TestesTodos            # 20 execucoes de cada
\end{Verbatim}
Números dos problemas: 1 produtor e consumidor, 2 barbeiro, 3 leitores e escritores, 4 fumantes, 5 filósofos,
6 selvagens, 7 três robôs.

\end{document}
